\chapter[Lezione 15 --- 28 ottobre 2026]{Lezione 15: limiti di funzioni monotone, esponenziali e logaritmi}
\label{chap:lezione-15}

\section{Limiti delle funzioni monotone}
\label{sec:lez15-limiti-monotone}

La monotonia introdotta nella \cref{def:lez10-monotonia}
permette di determinare i limiti agli estremi di un intervallo
mediante gli estremi dell'immagine della funzione.

\begin{theorem}[Limiti delle funzioni monotone]
\label{thm:lez15-limiti-monotone}
Siano \(a,b\in\overline{\R}\), con \(a<b\), e sia
\(I=(a,b)\subseteq\R\). Si consideri una funzione
\(f\colon I\longrightarrow\R\). Allora:
\begin{enumerate}
  \item se \(f\) è debolmente crescente su \(I\), si ha
  \[
    \lim_{x\to a^+}f(x)=\inf\{f(x):x\in I\},
    \qquad
    \lim_{x\to b^-}f(x)=\sup\{f(x):x\in I\};
  \]
  \item se \(f\) è debolmente decrescente su \(I\), si ha
  \[
    \lim_{x\to a^+}f(x)=\sup\{f(x):x\in I\},
    \qquad
    \lim_{x\to b^-}f(x)=\inf\{f(x):x\in I\}.
  \]
\end{enumerate}
Se \(a=-\infty\) o \(b=+\infty\), le rispettive scritture
\(x\to a^+\) e \(x\to b^-\) indicano i limiti per
\(x\to-\infty\) e \(x\to+\infty\).
Gli estremi dell'immagine sono intesi nella retta reale estesa.
\end{theorem}

È essenziale considerare l'immagine dell'intervallo aperto
\(I=(a,b)\). Se la funzione è definita anche in \(a\),
il suo valore in tale punto non influisce sul limite destro,
ma può modificare l'estremo inferiore dell'immagine.
Ad esempio, per una funzione debolmente crescente su
\(J=[a,b)\), può accadere che
\[
  f(a)=\inf f(J)
  <\inf\{f(x):x\in(a,b)\}
  =\lim_{x\to a^+}f(x),
\]
come illustrato nella \cref{fig:lez15-estremo-incluso}.

\begin{figure}[htbp]
  \centering
  \begin{tikzpicture}[unipd/diagram,x=1cm,y=.9cm]
    \draw[->,unipd-muted] (.3,-.05) -- (.3,4.05) node[above] {\(y\)};
    \draw[->,unipd-muted] (.1,0) -- (6.1,0) node[right] {\(x\)};
    \draw[densely dashed,unipd-muted] (.3,.7) -- (1.35,.7);
    \draw[densely dashed,unipd-muted] (.3,1.55) -- (1.35,1.55);
    \draw[very thick,unipd-primary]
      plot[smooth] coordinates {
        (1.35,1.55) (1.75,1.77) (2.35,1.91)
        (3.15,2.12) (3.9,2.35) (4.45,2.75) (5.45,3.35)
      };
    \draw[unipd-primary,fill=white,thick] (1.35,1.55) circle (2.5pt);
    \draw[unipd-primary,fill=white,thick] (5.45,3.35) circle (2.5pt);
    \fill[unipd-primary] (1.35,.7) circle (2.2pt);
    \node[below=6pt] at (1.35,0) {\(a\)};
    \node[below=6pt] at (5.45,0) {\(b\)};
    \node[left=5pt] at (.3,.7) {\(f(a)\)};
    \node[left=5pt] at (.3,1.55) {\(\inf f((a,b))\)};
  \end{tikzpicture}
  \caption{Una funzione debolmente crescente su \(J=[a,b)\)
  può assumere in \(a\) un valore inferiore al limite destro.
  L'estremo inferiore di \(f(J)\) non coincide allora con
  quello di \(f((a,b))\).}
  \label{fig:lez15-estremo-incluso}
\end{figure}

\begin{proof}
Si supponga che \(f\) sia debolmente crescente e che
\(b\in\R\). Si ponga
\[
  \lambda:=\sup\{f(x):x\in I\}\in\R\cup\{+\infty\}.
\]
Si dimostra che
\[
  \lim_{x\to b^-}f(x)=\lambda.
\]

\emph{Primo caso: \(\lambda\in\R\).}
Sia \(\varepsilon>0\). Per il
\cref{thm:caratterizzazione-sup}, esiste
\(x_\varepsilon\in I\) tale che
\[
  \lambda-\varepsilon<f(x_\varepsilon)\le\lambda.
\]
Per ogni \(x\in(x_\varepsilon,b)\), dalla monotonia di
\(f\) e dalla definizione di estremo superiore segue
\[
  \lambda-\varepsilon<f(x_\varepsilon)
  \le f(x)\le\lambda<\lambda+\varepsilon.
\]
Pertanto
\[
  |f(x)-\lambda|<\varepsilon
  \qquad\forall x\in(x_\varepsilon,b).
\]

\begin{figure}[H]
  \centering
  \begin{tikzpicture}[unipd/diagram,x=1cm,y=1cm]
    \draw[->,unipd-muted] (.1,2.15) -- (6.9,2.15)
      node[right] {\(f(x)\)};
    \draw[very thick,unipd-primary] (1.0,2.15) -- (5.9,2.15);
    \draw[unipd-muted] (1.0,2.03) -- (1.0,2.27);
    \draw[unipd-muted] (5.9,2.03) -- (5.9,2.27);
    \fill[unipd-primary] (2.7,2.15) circle (2pt);
    \fill[unipd-primary] (4.3,2.15) circle (2pt);
    \node[below=7pt] at (1.0,2.15) {\(\lambda-\varepsilon\)};
    \node[above=7pt] at (2.7,2.15) {\(f(x_\varepsilon)\)};
    \node[above=7pt] at (4.3,2.15) {\(f(x)\)};
    \node[below=7pt] at (5.9,2.15) {\(\lambda\)};

    \draw[->,unipd-muted] (.1,.4) -- (6.9,.4)
      node[right] {\(x\)};
    \draw[very thick,unipd-accent] (1.1,.4) -- (6.15,.4);
    \draw[unipd-accent,fill=white,thick] (1.1,.4) circle (2.4pt);
    \draw[unipd-accent,fill=white,thick] (6.15,.4) circle (2.4pt);
    \fill[unipd-accent] (3.7,.4) circle (2pt);
    \node[below=7pt] at (1.1,.4) {\(x_\varepsilon\)};
    \node[below=7pt] at (3.7,.4) {\(x\)};
    \node[below=7pt] at (6.15,.4) {\(b\)};
  \end{tikzpicture}
  \caption{Per \(x_\varepsilon<x<b\), la monotonia e la
  proprietà di estremo superiore forniscono
  \(\lambda-\varepsilon<f(x_\varepsilon)\le f(x)\le\lambda\).}
  \label{fig:lez15-supremo-finito}
\end{figure}

Ponendo
\[
  \delta_\varepsilon:=b-x_\varepsilon>0,
\]
si ottiene
\[
  |f(x)-\lambda|<\varepsilon
  \qquad\forall x\in(b-\delta_\varepsilon,b).
\]
Per la \cref{def:lez14-limiti-unilateri}, si conclude
\[
  \lim_{x\to b^-}f(x)=\lambda.
\]

\emph{Secondo caso: \(\lambda=+\infty\).}
L'immagine di \(f\) non è superiormente limitata.
Per ogni \(K\in\R\), esiste quindi \(x_K\in I\) tale che
\[
  f(x_K)>K.
\]
Per ogni \(x\in(x_K,b)\), dalla monotonia di \(f\)
segue
\[
  f(x)\ge f(x_K)>K.
\]

\begin{figure}[H]
  \centering
  \begin{tikzpicture}[unipd/diagram,x=1cm,y=1cm]
    \draw[->,unipd-muted] (.1,2.15) -- (6.9,2.15)
      node[right] {\(f(x)\)};
    \draw[very thick,unipd-primary] (1.0,2.15) -- (6.15,2.15);
    \draw[unipd-muted] (1.0,2.03) -- (1.0,2.27);
    \fill[unipd-primary] (2.7,2.15) circle (2pt);
    \fill[unipd-primary] (4.3,2.15) circle (2pt);
    \node[below=7pt] at (1.0,2.15) {\(K\)};
    \node[above=7pt] at (2.7,2.15) {\(f(x_K)\)};
    \node[above=7pt] at (4.3,2.15) {\(f(x)\)};

    \draw[->,unipd-muted] (.1,.4) -- (6.9,.4)
      node[right] {\(x\)};
    \draw[very thick,unipd-accent] (1.1,.4) -- (6.15,.4);
    \draw[unipd-accent,fill=white,thick] (1.1,.4) circle (2.4pt);
    \draw[unipd-accent,fill=white,thick] (6.15,.4) circle (2.4pt);
    \fill[unipd-accent] (3.7,.4) circle (2pt);
    \node[below=7pt] at (1.1,.4) {\(x_K\)};
    \node[below=7pt] at (3.7,.4) {\(x\)};
    \node[below=7pt] at (6.15,.4) {\(b\)};
  \end{tikzpicture}
  \caption{Se \(\lambda=+\infty\), per ogni \(K\in\R\)
  esiste \(x_K\in I\) tale che \(K<f(x_K)\le f(x)\)
  per ogni \(x\in(x_K,b)\).}
  \label{fig:lez15-supremo-infinito}
\end{figure}

Posto \(\delta_K:=b-x_K>0\), si ha dunque
\[
  f(x)>K
  \qquad\forall x\in(b-\delta_K,b).
\]
Per la \cref{def:lez14-limiti-unilateri}, si ottiene
\[
  \lim_{x\to b^-}f(x)=+\infty.
\]
\end{proof}

Il secondo punto del \cref{thm:lez15-limiti-monotone}
si riconduce al primo. Infatti, se \(f\) è debolmente
decrescente, allora \(-f\) è debolmente crescente
per la \cref{prop:lez10-monotonia-opposta}. Applicando
il primo punto a \(-f\), si ha
\[
\begin{aligned}
  \lim_{x\to a^+}(-f)(x)
    &=\inf\{-f(x):x\in I\},\\
  \lim_{x\to b^-}(-f)(x)
    &=\sup\{-f(x):x\in I\}.
\end{aligned}
\]
Per l'\cref{ex:estremi-insieme-opposto},
\[
  \inf\{-f(x):x\in I\}=-\sup\{f(x):x\in I\},
\]
\[
  \sup\{-f(x):x\in I\}=-\inf\{f(x):x\in I\}.
\]
Cambiando segno ai limiti mediante il
\cref{thm:lez14-limite-prodotto}, si ricavano le
due uguaglianze del secondo punto.

\section{Limiti di funzioni esponenziali e logaritmiche}
\label{sec:lez15-limiti-elementari}

\begin{figure}[htbp]
  \centering
  \begin{tikzpicture}
  \begin{axis}[
    width=0.88\textwidth,
    height=8cm,
    axis lines=middle,
    xmin=-4.2, xmax=4.2,
    ymin=-4.2, ymax=4.2,
    xtick={-4,-3,-2,-1,1,2,3,4},
    ytick={-4,-3,-2,-1,1,2,3,4},
    xlabel={\(x\)},
    ylabel={\(y\)},
    xlabel style={at={(axis description cs:1,0.48)},anchor=west},
    ylabel style={at={(axis description cs:0.52,1)},anchor=south},
    enlargelimits=false,
    clip=true,
    samples=250,
    domain=-4:4,
    legend style={
      at={(0.5,-0.16)},
      anchor=north,
      legend columns=2,
      draw=none,
      fill=none,
      font=\small
    },
    every axis plot/.append style={very thick}
  ]
    \addplot[unipd-primary] {2^x};
    \addlegendentry{\(y=2^x\)}

    \addplot[unipd-accent,domain=0.05:4] {ln(x)/ln(2)};
    \addlegendentry{\(y=\log_2 x\)}

    \addplot[teal!70!black] {(1/2)^x};
    \addlegendentry{\(y=\left(\frac12\right)^x\)}

    \addplot[violet!80!black,domain=0.05:4] {ln(x)/ln(1/2)};
    \addlegendentry{\(y=\log_{1/2}x\)}
  \end{axis}
\end{tikzpicture}

  \caption{Grafici rappresentativi delle funzioni esponenziali e logaritmiche
  nei casi \(a>1\) e \(0<a<1\). In figura sono scelte le basi
  \(a=2\) e \(a=\frac12\), che illustrano i comportamenti descritti
  nel \cref{thm:lez15-limiti-esponenziale-logaritmo}.}
  \label{fig:lez15-esponenziale-logaritmo}
\end{figure}

\clearpage

\begin{theorem}[Limiti dell'esponenziale e del logaritmo]
\label{thm:lez15-limiti-esponenziale-logaritmo}
Sia \(a>0\), con \(a\ne1\). Allora:
\begin{enumerate}
  \item se \(a>1\), si ha
  \[
    \lim_{x\to-\infty}a^x=0^+,
    \qquad
    \lim_{x\to+\infty}a^x=+\infty;
  \]
  \item se \(0<a<1\), si ha
  \[
    \lim_{x\to-\infty}a^x=+\infty,
    \qquad
    \lim_{x\to+\infty}a^x=0^+;
  \]
  \item se \(a>1\), si ha
  \[
    \lim_{x\to0^+}\log_a x=-\infty,
    \qquad
    \lim_{x\to+\infty}\log_a x=+\infty;
  \]
  \item se \(0<a<1\), si ha
  \[
    \lim_{x\to0^+}\log_a x=+\infty,
    \qquad
    \lim_{x\to+\infty}\log_a x=-\infty.
  \]
\end{enumerate}
La scrittura \(0^+\) indica che i valori della funzione
tendono a zero dall'alto, nel senso della
\cref{def:lez14-avvicinamento-valore}.
Per \(a=1\), escluso dalle ipotesi, si ha
\(a^x=1\) per ogni \(x\in\R\).
\end{theorem}

\begin{proof}
\emph{Primo punto.}
Sia \(a>1\). Per l'\cref{ex:lez10-logaritmo-monotono},
la funzione \(x\longmapsto a^x\) è strettamente crescente
su \(\R\). Applicando il \cref{thm:lez15-limiti-monotone}
all'intervallo \((-\infty,+\infty)\), si ha
\[
  \lim_{x\to-\infty}a^x
    =\inf\{a^x:x\in\R\},
  \qquad
  \lim_{x\to+\infty}a^x
    =\sup\{a^x:x\in\R\}.
\]

\emph{Limite per \(x\to+\infty\).}
Si ponga \(c:=a-1>0\). Per ogni \(n\in\N\),
la disuguaglianza di Bernoulli del
\cref{thm:disuguaglianza-bernoulli} fornisce
\[
  a^n=(1+c)^n\ge1+nc.
\]
Poiché \(1+nc\to+\infty\) per \(n\to+\infty\),
per il \cref{thm:lez12-confronto-secondo} si ottiene
\[
  \lim_{n\to+\infty}a^n=+\infty.
\]
L'insieme \(\{a^n:n\in\N\}\) non è quindi
superiormente limitato. Poiché
\[
  \{a^n:n\in\N\}\subseteq\{a^x:x\in\R\},
\]
neppure l'immagine dell'esponenziale è superiormente
limitata. Pertanto
\[
  \sup\{a^x:x\in\R\}=+\infty,
\]
e, per il \cref{thm:lez15-limiti-monotone},
\[
  \lim_{x\to+\infty}a^x=+\infty.
\]

\emph{Limite per \(x\to-\infty\).}
Dal limite appena dimostrato e dal
\cref{thm:lez14-limite-rapporto} segue
\[
  \lim_{n\to+\infty}a^{-n}
    =\lim_{n\to+\infty}\frac1{a^n}=0.
\]
Poiché \(a^{-n}>0\) per ogni \(n\in\N\), per ogni
\(\varepsilon>0\) esiste \(n\in\N\) tale che
\[
  0<a^{-n}<\varepsilon.
\]
Per il \cref{thm:caratterizzazione-inf}, si ha quindi
\[
  \inf\{a^{-n}:n\in\N\}=0.
\]
Dall'inclusione
\[
  \{a^{-n}:n\in\N\}\subseteq\{a^x:x\in\R\}
\]
segue
\[
  \inf\{a^x:x\in\R\}
  \le\inf\{a^{-n}:n\in\N\}=0.
\]
D'altra parte, \(a^x>0\) per ogni \(x\in\R\), quindi
\(0\) è un minorante dell'immagine e
\[
  0\le\inf\{a^x:x\in\R\}\le0.
\]
Si conclude che
\[
  \inf\{a^x:x\in\R\}=0,
  \qquad
  \lim_{x\to-\infty}a^x=0^+.
\]

\emph{Secondo punto.}
Sia \(0<a<1\) e si ponga
\[
  c:=\frac1a>1.
\]
Per il primo punto,
\[
  \lim_{x\to-\infty}c^x=0^+,
  \qquad
  \lim_{x\to+\infty}c^x=+\infty.
\]
Poiché \(a^x=1/c^x\), per il
\cref{thm:lez14-limite-rapporto} si ottiene
\[
  \lim_{x\to-\infty}a^x
    =\lim_{x\to-\infty}\frac1{c^x}
    =+\infty
\]
e
\[
  \lim_{x\to+\infty}a^x
    =\lim_{x\to+\infty}\frac1{c^x}
    =0^+.
\]

\emph{Terzo punto.}
Sia \(a>1\). La funzione
\[
  \varphi\colon\R\longrightarrow(0,+\infty),
  \qquad x\longmapsto a^x,
\]
è strettamente crescente. La sua applicazione inversa
\[
  \varphi^{-1}\colon(0,+\infty)\longrightarrow\R,
  \qquad x\longmapsto\log_a x,
\]
è quindi strettamente crescente per la
\cref{prop:lez10-monotonia-inversa}.

Poiché
\[
  \operatorname{Im}(\log_a)=\R,
\]
per il \cref{thm:lez15-limiti-monotone} si ha
\[
\begin{aligned}
  \lim_{x\to0^+}\log_a x
    &=\inf\{\log_a x:x>0\}
      =\inf\R=-\infty,\\
  \lim_{x\to+\infty}\log_a x
    &=\sup\{\log_a x:x>0\}
      =\sup\R=+\infty.
\end{aligned}
\]

\emph{Quarto punto.}
Sia \(0<a<1\). Posto \(c:=1/a>1\), si ha
\[
  a^x=\frac1{c^x}.
\]
Poiché \(c^x\) è strettamente crescente e positiva,
la funzione \(x\longmapsto a^x\) è strettamente
decrescente su \(\R\). Per la
\cref{prop:lez10-monotonia-inversa}, anche la sua
inversa \(x\longmapsto\log_a x\) è strettamente
decrescente su \((0,+\infty)\).

Poiché \(\operatorname{Im}(\log_a)=\R\), applicando
il secondo punto del \cref{thm:lez15-limiti-monotone}
si ottiene
\[
\begin{aligned}
  \lim_{x\to0^+}\log_a x
    &=\sup\{\log_a x:x>0\}
      =\sup\R=+\infty,\\
  \lim_{x\to+\infty}\log_a x
    &=\inf\{\log_a x:x>0\}
      =\inf\R=-\infty.
\end{aligned}
\]
\end{proof}
